% Generated draft: requires mathematical and source review.
\documentclass[11pt]{article}
\usepackage[margin=25mm]{geometry}
\usepackage{fontspec}
\setmainfont{Latin Modern Roman}
\usepackage{amsmath,amssymb,mathtools,unicode-math}
\setmathfont{Latin Modern Math}
\usepackage{enumitem,needspace,xurl,hyperref}
\setlength{\parindent}{0pt}
\setlength{\parskip}{5pt}
\setlength{\emergencystretch}{3em}
\newcommand{\N}{\mathbb{N}}
\newcommand{\Z}{\mathbb{Z}}
\newcommand{\Q}{\mathbb{Q}}
\newcommand{\R}{\mathbb{R}}
\newcommand{\C}{\mathbb{C}}
\newcommand{\F}{\mathbb{F}}
\newcommand{\A}{\mathbb{A}}
\newcommand{\PP}{\mathbb{P}}
\newcommand{\E}{\mathbb{E}}
\newcommand{\eps}{\varepsilon}
\newcommand{\dd}{\,\mathrm{d}}
\begin{document}
\section*{UnsolvedMath ID: 30001127}
\textbf{Source problem code:} catalogue entry
\subsection{Definitions and mathematical statement}
The supplied record concerns Liesegang precipitation patterns: precipitation bands in reaction--diffusion systems exhibiting a characteristic spatial scaling. It does not specify the reaction equations, spatial domain, initial data, boundary conditions, definition of a band, or scaling law. Consequently, it does not determine a single fully specified mathematical problem.

The following notation expresses the structure of the question conditionally, without supplying those missing choices. Once a reaction--diffusion family and its admissible data have been specified, let \(\mathcal X\) denote the space of admissible kinetic constants, diffusion constants, and boundary conditions, equipped with a specified topology. Let \(\mathcal P\subseteq\mathcal X\) consist of the data for which the corresponding system exhibits Liesegang precipitation patterns satisfying a specified spatial scaling criterion. Both the topology and the criterion defining \(\mathcal P\) remain unspecified in the source record; any dependence on initial data must also be made explicit.

The source question asks which patterns persist under perturbations and which open sets of kinetic constants, diffusion constants, and boundary conditions support that persistence. In the conditional notation above, its parameter-robustness component asks for a characterization of data \(\theta_0\in\mathcal P\) for which there exists an open set \(U\subseteq\mathcal X\) such that
\[
\theta_0\in U,
\qquad
\forall\theta\in U,\quad \theta\in\mathcal P.
\]
It also asks for the extent and nature of the robustness under perturbations. The admissible perturbations and the meaning of persistence must be specified before this broader component has a precise mathematical formulation. The displayed condition is only a conditional framework for the source question, not an assertion that such an open set exists or a complete definition of dynamical robustness.

The supplied dated literature assessment reports a rigorous mechanism establishing robust Liesegang patterns for an open class of reaction--diffusion systems, while retaining explicit open cases concerning general robustness and temporal dynamics. No equations, hypotheses, or precise remaining temporal-dynamics questions are supplied, so these cannot be stated as specific theorems or conjectures here.
\subsection{Short English statement}
Determine when Liesegang precipitation bands and their characteristic spatial scaling persist as reaction kinetics, diffusion constants, and boundary conditions vary. A precise formulation requires a specified system, domain, initial and boundary data, scaling criterion, and topology of perturbations. The supplied classification is partially solved; the stated partial results and remaining temporal-dynamics cases require independent review.
\subsection{Sources}
\begin{itemize}
\item Local UnsolvedMath registry, permanent ID \(30001127\), title ``Robustness of Liesegang Precipitation Patterns,'' source problem code \texttt{catalogue entry}, in the Oberwolfach Reports Open Problems collection. The record attributes the question to workshop open-problem discussions documented in Oberwolfach Reports; no individual proposer is supplied. The proposed year is recorded as 2008.
\item Supplied report citation: \emph{Dynamics of Patterns (2009)}, Oberwolfach Reports, DOI \texttt{10.4171/owr/2008/57}, pp.\ 3201--3271. \url{https://doi.org/10.4171/owr/2008/57}. This link is also supplied under the literature-source label ``original OWR report.'' The report title's year and the recorded proposed year are retained as supplied.
\item Supplied literature source, labeled ``robustness of Liesegang patterns'': \url{https://www-users.cse.umn.edu/~scheel/preprints/liesegang.pdf}. The record supplies no authors, publication details, theorem number, or exact hypotheses for this reference.
\item Literature assessment dated exactly \texttt{2026-08-21}: the supplied curation labels the problem ``partially solved'' and reports a rigorous reaction--diffusion mechanism proving robust Liesegang patterns for an open class of systems, with explicit open cases for general robustness and temporal dynamics. This is supplied dated curation, not an independently verified assessment; the record itself requests independent verification of time-sensitive status and cited claims.
\item The registry labels the statement \texttt{corrected\_verified} and reports that its mathematical content was checked against the official Oberwolfach report, with notation, context, and question form normalized without changing the source-level claim. This verification claim is supplied provenance and has not been independently checked in this draft.
\end{itemize}
\end{document}
